%!TEX root = /Users/markelikalderon/Documents/Git/formwithoutmatter/aristotle.tex
\chapter{Form Without Matter} % (fold)
\label{cha:form_without_matter}

\section{The Generalized Form of Empedoclean Puzzlement} % (fold)
\label{sec:the_generalized_form_of_empedoclean_puzzlement}

% section the_generalized_form_of_empedoclean_puzzlement (end)

\section{The Wax Analogy} % (fold)
\label{sec:the_wax_analogy}

The assimilation of sensible form without matter is explained in terms of an analogy with wax receiving the impression of a signet ring:
\begin{quote}
	Generally, about all perception, we can say that a sense is what has the power of receiving into itself the sensible forms of things without the matter, in the way in which a piece of wax takes on the impress of a signet-ring without the iron or gold; what produces the impression is a signet of bronze or gold, but not qua bronze or gold: in a similar way the sense is affected by what is coloured or flavoured or sounding not insofar as each is what it is, but insofar as it is of such and such a sort and according to its form. (\emph{De Anima} \textsc{ii}.12 424\( ^{a} \)18--23; Smith in \citealt[42--43]{Barnes:1984uq})
\end{quote}

The analogy is making two points. First the assimilation of sensible form is compared with the wax's reception of the impression sealed by the signet ring. Second, Aristotle emphasizes what the wax receives, namely, the form of the seal, it's shape. The wax receives the form imposed upon it by the signet ring, but if does not receive any of the matter that composes the ring, be it bronze or gold, say. So, by analogy, when an animal perceives the white of the sun, they assimilate the chromatic form of the sun but none of its matter.

\subsection{Assimilation of Form} % (fold)
\label{sub:assimilation_of_form}

The analogy is, of course, of Platonic provenance. In the \emph{Theaetetus}, Plato writes;
\begin{quotation}
	\textsc{socrates}: You will think better of it when you hear the rest. To judge truly is a fine thing and there is something discreditable in error.
	
	\textsc{theaetetus}: Of course.
	
	\textsc{socrates}: Well, they say the differences arise in this way. When a man has in his mind a good thick slab of wax, smooth and knealed to the right consistency, and the impressions that come through the senses are stamped on these tables of the `heart'---Homer's words hints at the mind's likeness to wax---then the imprints are clear and deep enough to last a long time. Such people are quick to learn and also have good memories, and besides they do not interchange the imprints of their perceptions but think truly. These imprints being distinct and well spaced are quickly assigned to their several stamps---the `real things' as they are called---and such men are said to be clever. Do you agree?
	
	\textsc{theaetetus}: Most emphatically.
	
	\textsc{socrates}: When a person has what the poet's wisdom commends as a `shaggy heart', or when the block is muddy or made of impure wax, or oversoft or hard, the people with soft wax are quick to learn, but forgetful, those with hard wax, the reverse. Where it is shaggy or rough, a gritty kind of stuff containing a lot of earth or dirt, the impressions obtained are indistinct; so are they too when the stuff is hard, for they have no depth. Impressions in soft wax also are indistinct, because they melt together and soon become blurred. And if besides this, they overlap through being crowded together into some wretched little narrow mind, they are still indistinct. All these types are likely to judge falsely. When they see or hear or think of something, they cannot quickly assign things to their several imprints. Because they are so slow and sort things into their wrong places, they constantly see and hear and think amiss, and say they are mistaken about things and stupid. (Plato, \emph{Theaetetus} 194\( ^{c} \)--195\( ^{a} \); Cornford in \citealt{Hamilton:1961fk})
\end{quotation}

One salient difference between the Aristotelian and Platonic analogies is that whereas Aristotle deploys the wax analogy to explain perception, Plato does so to explain judgment. This difference is most likely intentional and pointed. As we discussed in chapters~\ref{sec:definition} and \ref{sec:the_objects_of_perception}, and as \citet{Sorabji:1971fr,Sorabji:2003fk} emphasizes, Aristotle is extending the domain of perception as Plato conceives of it. Not only are the objects of perception no longer confined to the primary objects---we perceive common sensibles as well---but Aristotle also maintains that we can discriminate among sensory objects and this is the exercise of our perceptual capacities. Plato, in contrast, maintained that what is ``common'' to the objects of sense---that they are each the same and different from the others---is determined by cognitive, not perceptual capacities. Aristotle underscores this extension of our perceptual capacities by his use of the Platonic analogy. Aristotle emphasizes the fact that he is assigning to perception some of the functions that Plato assigns to judgment by using the analogy that Plato used to explain judgment to explain perception instead.

Another, perhaps less salient, difference between the Aristotelian and Platonic analogies concerns the nature of the agent acting upon the wax. Whereas what makes an impression on the wax for Plato is a stylus, what makes an impression on the wax for Aristotle is a signet ring. Plato has in mind a wax tablet, used for writing. Plato thus belongs to the tradition, and perhaps inaugurates it, of using the then current writing technology as a model for the mind. Think of Locke's blank slate, or the functionalist slogan that the mind is the software of the brain, itself coinciding with the emergence of text-editing and word-processing. Perhaps Nietzsche was right that our writing tools act upon our thoughts---they at least have a tendency to influence our philosophy of mind. Aristotle varies this aspect of the Platonic analogy. It is not a stylus on wax that creates the impression, it is a signet ring. Why the substitution? I believe the variation is intentional. It's significance will emerge in subsequent discussion.

Not only does Aristotle deploy the Platonic analogy to explain his conception of perception, but he also transforms how the analogy is understood. Plato's explanation of the reliability of memory and judgment crucially relies on causal features of the situation. An objects' impression is the effect it has on the mind's wax. Importantly, however, Aristotle has in mind a non-causal sense of impression.

To get a sense of this, first consider how Hume himself appropriates the Platonic analogy:
\begin{quote}
	All the perceptions of the human mind revolves themselves into two distinct kinds, which I shall call \textsc{impressions} and \textsc{ideas}. The difference betwixt these consists in the degree of force or liveliness, with which they strike upon the mind, and make their way into thought and consciousness. (Hume, \emph{Treatise} \textsc{i}.1.1.1)
\end{quote}
Hume begs his reader's indulgence in taking liberty with the use of the terms ``impression'' and ``idea'' (\emph{Treatise} \textsc{i}.1.1 n2) though he proclaims that it is at least a virtue of his regimentation that he ``restore[s] the word, \emph{idea}, to its original sense, from which Mr. \emph{Locke} had perverted it, in making it stand for all our perceptions.'' Hume is presumably taking liberty with his use of the term ``perception'' as well. Locke perverts the use of ``idea'' by making it stand for all perceptions. But perception, here, is not confined to the objects of sensory experience. By ``perception'' Hume means whatever is or could be the object of the mind (admittedly in a post-Cartesian conception of mind unavailable to the ancients). The perceptions present to the human mind are either impressions or ideas, depending upon the force or liveliness with which they strike the mind. Impressions are the objects that are presented with ``the most force and violence'' (\emph{Treatise} \textsc{i}.1.1.1).  Concerning impressions, Hume further explains:
\begin{quote}
	By the term \emph{impression} I wou'd not be understood to express the manner, in which our lively perceptions are produc'd in the soul, but merely the perceptions themselves; for which there is no particular name either in \emph{English} or any other language, that I know of. (\emph{Treatise}, \textsc{i}.1.1 n2)
\end{quote}
Consistent with this qualification, Hume is, nevertheless, operating with a causal notion of impression. The lively perceptions presented before the mind in viewing the streets of Edinburgh are themselves the effect of external causes. Inquiring after the manner in which such lively perceptions are produced in the soul does not, however, fall within the purview of Hume's new science of human nature. It belongs, rather, to a particular branch of natural philosophy, speculative anatomy. The qualification does not deny that impressions are effects. It rather signals Hume's intention to confine himself to, what might be described in a later terminology as, the phenomenology, the intentional objects of consciousness, in particular, to the few regular principles that govern the persistent change among the perceptions presented to the mind.

Like Aristotle, Hume is departing from Plato in using the analogy of impressions on wax to describe perceptual experience (and more besides, passions are impressions as well). But Hume retains a feature of Plato's original use of the analogy. Hume is still thinking of sensory impressions as the effects produced in the perceiver by external objects acting upon them. Is there an intelligible alternative? How else might talk of impressions be understood? 

Consider the closely related metaphor of shaping. There is clearly a causal sense of shaping. When the signet ring shapes the wax it causes the wax to be modified in a certain way. The wax takes on the shape imposed upon it by the signet ring. Similarly, Nazi bombing shaped the London skyline. It caused that skyline to be configured in a certain way, the way imposed upon it by the bombing. Importantly, however, there is another sense of shaping, not a causal sense, but a constitutive sense. Whereas Nazi bombing shaped the London skyline merely in a causal sense, St Paul's constitutively shapes that skyline by being a contour of it. This is dramatically demonstrated in Herbert Mason's iconic photograph (see figure~\ref{fig:stpauls}). St Paul's shapes the London skyline by being a part of it, despite the causal impact of Nazi bombing.

\begin{figure}[htbp]
	\centering
		\includegraphics[scale=0.6]{graphics/stpauls.jpg}
	\caption{St Paul's 29 December 1940}
	\label{fig:stpauls}
\end{figure}

Humean sensory impressions are shaped by the environment merely in a causal sense. This is central to Hume's use of the Platonic analogy. Just as a signet ring impinging upon the wax causes an impression, the environment impinging upon a perceiver with the appropriate sensory capacities causes a sensory impression. How exactly such sensory impressions are produced is a matter for the speculative anatomist. Hume's new science of human nature confines itself to sensory impressions and the regularities that can be discerned in the flux of sensory experience. But sensory impressions, episodes in the sensory flux, remain effects, nonetheless. But perhaps perceptual sensitivity is more than the environment impinging upon the state of a conscious subject. Perhaps there is more to perception than objects eliciting a conscious modification of the perceiving subject. Perhaps the environment can shape sensory consciousness in a constitutive, rather than merely a causal, sense.

Before exploring this idea further, let us first consider another important aspect of the Platonic analogy. Caston observes that the impression produced by a signet ring is linked to that particular ring and, hence, metonymically at least, to the legitimate possessor of that ring:
\begin{quote}
	A signet produces a sealing, an impression that establishes the identity of its owner and consequently his authority, rights, and prerogatives. When a sealing is placed on a document, especially for legal or official use, it authorizes the claims, obligations, promises, or orders made therein. A sealing thus differs from other impressions in that it \emph{purports to originate from a particular signet}. \citep[302]{Caston:2005cr}
\end{quote}
The impression of a signet ring plays a similar role to signatures. Just as a signature is linked to the particular person whose signature it is, the impression sealed upon the wax by a signet ring is linked to the legitimate possessor of the ring. Of course, signatures can be forged, as can signet rings, which can also be stolen, but these practices gain there point precisely by the link between a signature and sealed impression, on the one hand, and their legitimate possessors, on the other. 

Taking this feature of the Platonic analogy seriously has an important consequence for how sensory impressions are individuated. Just as a forged signature is not my signature, an impression sealed by a forged ring, or by a stolen ring, is not the seal of the ring's legitimate possessor. Impressions are individuated by their legitimate sources. If this feature of the analogy carries over, then perceptions, conceived on the model of sealed impressions, are individuated by their objects. A perception of Castor and a perception of Pollux are different perceptions, no matter how closely the twins may resemble one another.

Notice that understanding sensory impressions on the model of the effects of causal shaping does not have this consequence. If, as Hume maintained, cause and effect are contingently connected, the same effect, the same impression, could have been produced by a different cause. Sensory impressions, understood as the effects of causal shaping, are not individuated by their causes. If sensory impressions are individuated by their objects, they cannot be understood on the model of effects of causal shaping. How else might they be understood?

If sensory impressions are individuated by their objects, perhaps these objects shape sensory consciousness not causally, or at least not merely. Perhaps sensory impressions are individuated by their objects because their objects constitutively shape sensory consciousness.



% subsection assimilation_of_form (end)

\subsection{Without Matter} % (fold)
\label{sub:without_matter}



% subsection without_matter (end)

% section the_wax_analogy (end)

\section{The Resolution of Empedoclean Puzzlement} % (fold)
\label{sec:the_resolution_of_empedoclean_puzzlement}



% section the_resolution_of_empedoclean_puzzlement (end)

\section{Presence and Modernity} % (fold)
\label{sec:presence_and_modernity}

If the metaphysics of presence cannot be coherently sustained within modern philosophy, this is not because of any contradiction or incoherence within the metaphysics of presence, but because modern philosophy abandoned presence from the beginning.

% section presence_and_modernity (end)

% chapter form_without_matter (end)
